\clearpage
\section{Reading Indices}
\label{app:reading-indices}
\index{readability indices}

The readability analysis supplied for the Draft~1 plagiarism-check text reported the grade-level estimates shown in Table~\ref{tab:reading-indices}. The values are reproduced exactly as displayed by the analysis tool.

\begin{table}[htbp]
    \centering
    \caption{Reading indices reported for Draft~1}
    \label{tab:reading-indices}
    \begin{tabular}{@{}lr@{}}
        \toprule
        \textbf{Reading index} & \textbf{Grade-level estimate} \\
        \midrule
        Gunning Fog & 17.99 \\
        Flesch--Kincaid Grade Level & 15.26 \\
        SMOG & 15.84 \\
        Coleman--Liau & 19.82 \\
        Automated Readability Index & 17.92 \\
        \midrule
        Average grade level & 17.37 \\
        Median grade level & 17.92 \\
        \bottomrule
    \end{tabular}
\end{table}

Taken together, these estimates indicate that the current report uses language associated with advanced undergraduate to graduate or professional reading levels. This is plausible for a technical computer-architecture report, in which specialised terminology and long technical expressions increase formula-based scores. The indices estimate linguistic complexity only; they do not measure technical correctness, writing quality, originality, or the strength of the argument. The interpretation follows the general description of readability measures provided by the analysis site.\footnote{\url{https://www.analyzemywriting.com/readability_indices.html}}

\begin{figure}[htbp]
    \centering
    \includegraphics[width=0.98\textwidth]{"../Draft 1/appendices/assets/readability-indices-2026-08-24.png"}
    \caption{Readability-index results supplied for the Draft~1 text on 24 August 2026.}
    \label{fig:reading-indices-result}
\end{figure}
